% GENERATED FILE. Edit canonical lessons, then run npm run generate:books.
% canonical-source-hash: fnv1a64:1bc675f6ea2f2425
% canonical-lessons: KA-C221-sebu, KA-C221-kittale, KA-C221-pappayi, KA-C221-kallangadi, KA-C221-capati

\chapter{Apple, Orange, Papaya, Watermelon, Chapati}
\label{ch:ka-apple-orange-papaya-watermelon-chapati}
\hlchaptermodality{221}

\begin{chapteropening}
\textbf{By the end of this chapter:} \emph{I can say an apple, an orange, a papaya, a watermelon and a chapati.}

Complete the last lesson of chapter 221: I can say an apple, an orange, a papaya, a watermelon and a chapati.
\end{chapteropening}

\section[\texorpdfstring{\kn{ಸೇಬು} (sēbu)}{ಸೇಬು (sēbu)}]{\kn{ಸೇಬು} (sēbu) — an apple}

\label{lesson:KA-C221-sebu}



\begin{quote}
Before the new one: say the Kannada for a fan, then the Kannada for rubbish.
\end{quote}

\subsection*{You'll want to know: \kn{ಸೇಬು}}
\textbf{\kn{ಸೇಬು}} — \emph{sēbu} — \textquotedblleft{}an apple\textquotedblright{}.

\begin{grammarlens}[title={What you've built}]
One more word for this chapter.
\end{grammarlens}

\subsection*{Your turn}
\begin{itemize}
\raggedright
  \item \emph{Say it:} \emph{sēbu}
  \item \emph{Say it:} \emph{sēbu}, once more
  \item \emph{Say it:} say \emph{sēbu}
  \item \emph{Recall:} say the Kannada for a fan, then the Kannada for rubbish, then say \emph{sēbu} again
\end{itemize}

\begin{tcolorbox}[breakable,colback=teal!4,colframe=teal!35!black,title={Before you move on}]
What does \kn{ಸೇಬು} mean? (\textquotedblleft{}An apple\textquotedblright{}.) Say it once more.
\end{tcolorbox}
\section[\texorpdfstring{\kn{ಕಿತ್ತಳೆ} (kittaḷe)}{ಕಿತ್ತಳೆ (kittaḷe)}]{\kn{ಕಿತ್ತಳೆ} (kittaḷe) — an orange}

\label{lesson:KA-C221-kittale}



\begin{quote}
Before the new one: say the Kannada for rubbish, then the Kannada for an apple.
\end{quote}

\subsection*{You'll want to know: \kn{ಕಿತ್ತಳೆ}}
\textbf{\kn{ಕಿತ್ತಳೆ}} — \emph{kittaḷe} — \textquotedblleft{}an orange\textquotedblright{}.

\begin{grammarlens}[title={What you've built}]
One more word for this chapter.
\end{grammarlens}

\subsection*{Your turn}
\begin{itemize}
\raggedright
  \item \emph{Say it:} \emph{kittaḷe}
  \item \emph{Say it:} \emph{kittaḷe}, once more
  \item \emph{Say it:} say \emph{kittaḷe}
  \item \emph{Recall:} say the Kannada for rubbish, then the Kannada for an apple, then say \emph{kittaḷe} again
\end{itemize}

\begin{tcolorbox}[breakable,colback=teal!4,colframe=teal!35!black,title={Before you move on}]
What does \kn{ಕಿತ್ತಳೆ} mean? (\textquotedblleft{}An orange\textquotedblright{}.) Say it once more.
\end{tcolorbox}
\section[\texorpdfstring{\kn{ಪಪ್ಪಾಯಿ} (pappāyi)}{ಪಪ್ಪಾಯಿ (pappāyi)}]{\kn{ಪಪ್ಪಾಯಿ} (pappāyi) — a papaya}

\label{lesson:KA-C221-pappayi}



\begin{quote}
Before the new one: say the Kannada for an apple, then the Kannada for an orange.
\end{quote}

\subsection*{You'll want to know: \kn{ಪಪ್ಪಾಯಿ}}
\textbf{\kn{ಪಪ್ಪಾಯಿ}} — \emph{pappāyi} — \textquotedblleft{}a papaya\textquotedblright{}.

\begin{grammarlens}[title={What you've built}]
One more word for this chapter.
\end{grammarlens}

\subsection*{Your turn}
\begin{itemize}
\raggedright
  \item \emph{Say it:} \emph{pappāyi}
  \item \emph{Say it:} \emph{pappāyi}, once more
  \item \emph{Say it:} say \emph{pappāyi}
  \item \emph{Recall:} say the Kannada for an apple, then the Kannada for an orange, then say \emph{pappāyi} again
\end{itemize}

\begin{tcolorbox}[breakable,colback=teal!4,colframe=teal!35!black,title={Before you move on}]
What does \kn{ಪಪ್ಪಾಯಿ} mean? (\textquotedblleft{}A papaya\textquotedblright{}.) Say it once more.
\end{tcolorbox}
\section[\texorpdfstring{\kn{ಕಲ್ಲಂಗಡಿ} (kallaṅgaḍi)}{ಕಲ್ಲಂಗಡಿ (kallaṅgaḍi)}]{\kn{ಕಲ್ಲಂಗಡಿ} (kallaṅgaḍi) — a watermelon}

\label{lesson:KA-C221-kallangadi}



\begin{quote}
Before the new one: say the Kannada for an orange, then the Kannada for a papaya.
\end{quote}

\subsection*{You'll want to know: \kn{ಕಲ್ಲಂಗಡಿ}}
\textbf{\kn{ಕಲ್ಲಂಗಡಿ}} — \emph{kallaṅgaḍi} — \textquotedblleft{}a watermelon\textquotedblright{}.

\begin{grammarlens}[title={What you've built}]
One more word for this chapter.
\end{grammarlens}

\subsection*{Your turn}
\begin{itemize}
\raggedright
  \item \emph{Say it:} \emph{kallaṅgaḍi}
  \item \emph{Say it:} \emph{kallaṅgaḍi}, once more
  \item \emph{Say it:} say \emph{kallaṅgaḍi}
  \item \emph{Recall:} say the Kannada for an orange, then the Kannada for a papaya, then say \emph{kallaṅgaḍi} again
\end{itemize}

\begin{tcolorbox}[breakable,colback=teal!4,colframe=teal!35!black,title={Before you move on}]
What does \kn{ಕಲ್ಲಂಗಡಿ} mean? (\textquotedblleft{}A watermelon\textquotedblright{}.) Say it once more.
\end{tcolorbox}
\section[\texorpdfstring{\kn{ಚಪಾತಿ} (capāti)}{ಚಪಾತಿ (capāti)}]{\kn{ಚಪಾತಿ} (capāti) — a chapati}

\label{lesson:KA-C221-capati}



\begin{quote}
Before the new one: say the Kannada for a papaya, then the Kannada for a watermelon.
\end{quote}

\subsection*{You'll want to know: \kn{ಚಪಾತಿ}}
\textbf{\kn{ಚಪಾತಿ}} — \emph{capāti} — \textquotedblleft{}a chapati\textquotedblright{}.

\begin{grammarlens}[title={What you've built}]
One more word for this chapter.
\end{grammarlens}

\subsection*{Your turn}
\begin{itemize}
\raggedright
  \item \emph{Say it:} \emph{capāti}
  \item \emph{Say it:} \emph{capāti}, once more
  \item \emph{Say it:} say \emph{capāti}
  \item \emph{Recall:} say the Kannada for a papaya, then the Kannada for a watermelon, then say \emph{capāti} again
\end{itemize}

\begin{tcolorbox}[breakable,colback=teal!4,colframe=teal!35!black,title={Before you move on}]
What does \kn{ಚಪಾತಿ} mean? (\textquotedblleft{}A chapati\textquotedblright{}.) Say it once more.
\end{tcolorbox}
